% ============================================================
\section{Perguntas de pesquisa e hipóteses}\label{sec:perguntas}
% ============================================================

O objetivo geral, como na proposta, é caracterizar por meio de redes complexas quanto e de que
forma a estrutura de similaridade entre representações de tokens se transforma entre o espaço
lexical de entrada e as representações contextuais das camadas inicial, intermediária e final
de um Transformer. A hipótese central \emph{não} é que a identidade lexical desapareça, mas que
exista uma reorganização progressiva das vizinhanças e das comunidades, cuja intensidade possa
ser medida. O objetivo se desdobra em quatro perguntas.

\begin{description}
    \item[P1. Quanto e de que forma a organização da rede muda?]
    Quantificar a reorganização \emph{local}, pela similaridade de Jaccard entre as vizinhanças
    de uma mesma ocorrência em duas redes e pela dominância lexical (fração de vizinhos do mesmo
    tipo), e a reorganização \emph{global}, pela distribuição de graus de entrada (surgimento ou
    desaparecimento de \emph{hubs}), pela reciprocidade, pela clusterização, pelas distâncias e
    pelas comunidades.

    \emph{Hipótese H1.} Na rede lexical por ocorrência a dominância é máxima por construção
    (ocorrências do mesmo tipo são idênticas). Ela cai ao longo das camadas, mas continua acima
    do acaso no bloco 36. A mudança não é só local: esperamos \emph{hubs} diferentes em cada
    camada (em espaços de alta dimensão, alguns pontos viram vizinhos de muitos
    outros~\cite{radovanovic2010}) e comunidades menos alinhadas ao tipo.

    \item[P2. As mudanças concentram-se em determinadas camadas?]
    Comparar as transições lexical $\to$ bloco 1, bloco 1 $\to$ 18 e bloco 18 $\to$ 36, e as
    acumuladas, identificando a que concentra a maior queda de Jaccard e de dominância lexical.
    A curva camada a camada (os 36 blocos) localiza a mudança com mais precisão. A resposta é
    dada por faixa de frequência e por categoria de token (palavra inteira ou fragmento).

    \emph{Hipótese H2.} Fragmentos de palavras (início e continuação) precisam do contexto
    imediato para ``virar'' uma palavra; esperamos que a vizinhança deles mude mais cedo (já no
    bloco 1) do que a de palavras inteiras de conteúdo. Tokens frequentes, cujo
    \emph{embedding} é pouco informativo sozinho (palavras funcionais), também devem mudar
    cedo.

    \item[P3. Surgem comunidades associadas a contextos semelhantes?]
    Verificar se as comunidades deixam de ser explicadas pelo tipo de token e passam a ser
    explicadas pelo contexto de origem (tema, artigo, parágrafo, sentença) ou pelo que vem
    depois (o próximo token), com medidas de concordância entre partições (NMI, ARI, pureza)
    corrigidas por uma linha de base de permutação, e por inspeção qualitativa.

    \emph{Hipótese H3.} A concordância com o tipo cai e a concordância com o contexto sobe ao
    longo das camadas. No bloco 36, a concordância com o \emph{próximo} token deve crescer,
    porque é essa a informação que a última camada precisa carregar. A faixa de posição entra
    como controle: se ela explicar as comunidades, o resultado é um artefato posicional, e não
    contexto.

    \item[P4. Ocorrências de um mesmo token separam-se conforme o contexto?]
    Para tipos com muitas ocorrências ($f_t\ge 10$), em especial palavras usadas em sentidos
    diferentes em temas diferentes (as \emph{palavras-alvo}, como \emph{banco}, \emph{campo} e
    \emph{órgão}), verificar se representações inicialmente idênticas passam a ocupar
    comunidades ou regiões diferentes da rede.

    \emph{Hipótese H4.} As palavras-alvo se separam mais, por tema de sentido, do que as
    palavras de controle (\emph{ano}, \emph{grande}, \emph{água}\ldots), que têm frequência e
    distribuição por tema comparáveis mas um sentido estável.
\end{description}

\subsection{Como cada pergunta será respondida}

A Tabela~\ref{tab:perguntas} liga cada pergunta às redes, medidas e seções de resultados. Duas
regras valem para todas as respostas:
\begin{itemize}
    \item \textbf{Nada é interpretado abaixo do nível de ruído.} Na rede lexical por
    ocorrência, a escolha entre vizinhos empatados é arbitrária. Uma mudança lexical $\to$
    contextual só conta como reorganização se for maior que a variação produzida pelo próprio
    desempate (o \emph{piso de ruído}, Seção~\ref{sec:met-ruido}). Da mesma forma, NMI só é
    lido depois de descontada a linha de base de permutação.
    \item \textbf{Resultados sempre estratificados.} Faixa de frequência, categoria de token,
    estrato da amostra e faixa de posição mudam os valores esperados das medidas (por exemplo, a
    dominância lexical máxima depende de $f_t$). Por isso as distribuições são reportadas por
    grupo, e não só como médias globais.
\end{itemize}

\begin{table}[htbp]
\centering
\caption{Perguntas, redes usadas, medidas e onde os resultados aparecem.}\label{tab:perguntas}
\small
\begin{tabularx}{\linewidth}{@{}l >{\raggedright\arraybackslash}p{3.6cm} X l@{}}
\toprule
& \textbf{Redes} & \textbf{Medidas} & \textbf{Seção} \\
\midrule
P1 & por ocorrência (lexical, blocos 1, 18, 36), direcionada e por união &
$J_i$ e $D_i$ por faixa, categoria e estrato; CCDF do grau de entrada, \emph{hubs},
reciprocidade; métricas globais; comunidades & \ref{sec:res-p1} \\
P2 & as mesmas, mais as 36 camadas &
$1-J$ acima do piso de ruído e $\Delta D$ por transição; fração de vértices cuja maior mudança
cai em cada transição; curvas $J(\ell,\ell+1)$ e $D(\ell)$ & \ref{sec:res-p2} \\
P3 & só o núcleo da amostra, por união &
NMI (bruto e menos a linha de base), ARI e pureza contra tipo, tema, artigo, parágrafo,
sentença, próximo token e faixa de posição, camada a camada & \ref{sec:res-p3} \\
P4 & rede inteira; tipos com $f_t\ge10$ &
número efetivo de comunidades $e^{H}$, NMI comunidade$\times$tema dentro do tipo,
autossimilaridade, cosseno intra $-$ entre temas, separação por sentido anotado, redes ego &
\ref{sec:res-p4} \\
Vocabulário & rede lexical de tipos (151,6 mil vértices) &
graus, \emph{hubs}, clusterização, distâncias; comunidades contra a classe de escrita;
vizinhança das palavras-alvo & \ref{sec:res-vocab} \\
\bottomrule
\end{tabularx}
\end{table}
